\documentclass[10pt,a4paper]{article}
\usepackage[utf8]{inputenc}
\usepackage[a4paper,left=3cm,right=2cm]{geometry}
\usepackage{amsmath}
\usepackage[thmmarks, amsmath, thref]{ntheorem}
\usepackage{amsfonts}
\usepackage{amssymb}
\usepackage{fancyhdr}
\pagestyle{fancy}
\usepackage{anysize}
\usepackage{graphicx}
\usepackage{float}
\usepackage{hyperref}
\usepackage{multirow}
\usepackage{enumitem}
\newtheorem{theorem}{Theorem}
\theoremsymbol{\ensuremath{\blacksquare}}
\newtheorem*{solution}{Solución:}

\usepackage{xcolor}

% Margenes y formalidades

\textwidth 6.25in % Margen Derecho y ancho del texto
\textheight 8.5in
\oddsidemargin 0in % Margen Izquierdo
\headheight 0in % Largo del texto
\leftmargin 1in
\parindent 0em % Salto entre la indentación (Sangría)https://www.overleaf.com/4519127669sdpptkhvsbfn
\topmargin -0.5in % Margen Superior
\parskip 0.5ex % Salto entre parrafos
\baselineskip 1pt

%\lhead{\includegraphics[scale=0.2]{logo.png}} % texto izquierda de la cabecera
%\chead{}                                      % texto centro de la cabecera
%\rhead{\includegraphics[scale=0.2]{FIG/logo-mba-2016.png}} % número de página a la derecha

\renewcommand{\headrulewidth}{0.4pt} % grosor de la línea de la cabecera
\renewcommand{\footrulewidth}{0.4pt} % grosor de la línea del pie

\begin{document}



\pagebreak
\vspace*{4mm} 

\begin{center}
\begin{huge}
\textbf{\\Ejercicios Econometría 1}
\end{Large}\\
%{\large 2023-2}\\
\end{center}


\section{Muestreo}

\textbf{Ejercicio 1:} Se tiene una población de 1000 individuos con una media de ingreso mensual de $1500$ y una desviación estándar de $300$. Si se toma una muestra aleatoria simple de 100 individuos, ¿cuál es la media esperada y la desviación estándar de la media muestral?



\section{Estadística Descriptiva}

\textbf{Ejercicio 2:} Se tiene el siguiente conjunto de datos sobre el número de ventas realizadas por 5 vendedores en un mes:
$$X = \{10,15,12,18,20\}$$
Calcule la media, mediana y varianza muestral.





\section{Regresión Lineal Simple}

\textbf{Ejercicio 3:} Se tiene la siguiente muestra de datos:
\begin{center}
\begin{tabular}{|c|c|}
\hline
X (Publicidad) & Y (Ventas) \\
\hline
2 & 4 \\
\hline
4 & 7 \\
\hline
6 & 10 \\
\hline
8 & 13 \\
\hline
10 & 16 \\
\hline
\end{tabular}
\end{center}

Ajuste un modelo de regresión lineal simple de la forma:
$$Y = \beta_0 + \beta_1 X + \epsilon $$

usando el método de mínimos cuadrados ordinarios.



\section{Regresión Lineal Simple: Educación y Salario}

\textbf{Ejercicio:} Se tiene la siguiente muestra de datos sobre los años de educación y el salario mensual (en dólares) de 10 individuos:

\begin{center}
\begin{tabular}{|c|c|}
\hline
Años de Educación ($X$) & Salario Mensual ($Y$) \\
\hline
12  & 600  \\
\hline
13  & 750  \\
\hline
14  & 900  \\
\hline
15  & 1100 \\
\hline
16  & 1300 \\
\hline
18  & 1600 \\
\hline
20  & 1850 \\
\hline
22  & 2100 \\
\hline
24  & 2300 \\
\hline
25  & 2500 \\
\hline
\end{tabular}
\end{center}

Ajuste un modelo de regresión lineal simple de la forma:
$$Y = \beta_0 + \beta_1 X + \epsilon $$


usando el método de mínimos cuadrados ordinarios, con el cálculo explícito de la varianza y la covarianza mediante tablas.


\section{Regresión Lineal Simple: Descomposición de Varianza y $R^2$}

\textbf{Ejercicio:} Se tiene la siguiente muestra de datos sobre la experiencia laboral (en años) y el ingreso mensual (en dólares) de 7 individuos:

\begin{center}
\begin{tabular}{|c|c|}
\hline
Experiencia Laboral ($X$) & Ingreso Mensual ($Y$) \\
\hline
2   & 1400  \\
\hline
4   & 1800  \\
\hline
6   & 2200  \\
\hline
8   & 2600  \\
\hline
10  & 3100  \\
\hline
12  & 3600  \\
\hline
14  & 4200  \\
\hline
\end{tabular}
\end{center}

Ajuste un modelo de regresión lineal simple de la forma:

$$Y = \beta_0 + \beta_1 X + \epsilon $$

usando el método de mínimos cuadrados ordinarios y calcule la descomposición de la varianza: SST, SSE y SSR, junto con el coeficiente de determinación $R^2$.



\end{document}